\scanpage{158}
\mspage{103}
Cela implique donc que
\[
C(X)\xrightarrow[\mathrm{pr}_1\pi']{}X
\]
est $\in W$ (étant déduit de $\sigma_S:\mathrm{Ch}(S)\longrightarrow S$ par chgt. de base), donc \deleted{toute} section \inserted{$i_X:X\longrightarrow C(X)$} de $C(X)$ sur $X$ est dans $i_X$. On a donc bien une factorisation, fonctorielle en $X$
\margin{\unsure{''}}
\diagram{\begin{tikzcd}[ampersand replacement=\&]
{} \& C(X)\arrow[dr,"{\overline f_X=f_{C(X)}}"] \& {}\\
X\arrow[ur,"i_X"]\arrow[rr,"{f_X\deleted{\unsure{X}}}"] \& {} \& S
\end{tikzcd}}
avec
\[
i_X\in W,\quad f_{C(X)}\in\mathrm{Fib}\text{.}
\]
C'est tout ce qu'on avait utilisé pour établir
\[
(\underline{\mathrm{Fib}}\,S)W_S^{-1}\mathrel{\overset{\approx}{\rule{1.7em}{.4pt}}}(\mathcal M_{/S})W_S^{-1}\text{.}
\]

Mais il s'agit maintenant, sous l'hyp. où on introduit un $\mathcal L$, de prouver que pour tt $X$ dans $\mathcal M_{/S}$, les deux flèches
\diagram{\begin{tikzcd}[ampersand replacement=\&]
C(X)\arrow[r,shift left=.6ex,"C(i_X)"]\arrow[r,shift right=.6ex,"i_{C(X)}"'] \& C(C(X))
\end{tikzcd}}
sont égales dans\unsure{.} $(\mathcal M_{/S})(W_S^{\mathcal L})^{-1}$. Comme par hyp. on a $W_S^{\mathcal L}=W_S$ dans $\underline{\unsure{\mathrm{Fb}|\mathrm{Fib}}}\,S$ (Hyp 1), et que $\deleted{\unsure{B}}C(X),C(C(X))\in\mathrm{Ob}\,\underline{\mathrm{Fib}}\,S$,
% archive sheet number: 157
